\setcounter{table}{7}
\begin{table}[H]
\centering
\caption{Implications for apparent two-wave transitions}
\scriptsize
\setlength{\tabcolsep}{7pt}
\begin{tabular}{@{}lrrr@{}}
\toprule
Model & Classification only & Genuine reversal & Genuine persistence \\
\midrule
\multicolumn{4}{l}{\textit{All apparent transitions}} \\
Table 1 (1): raw/no error & 0.0 & 8.4 & 91.6 \\
\addlinespace[1pt]
Table 1 (2): symmetric error & 67.0 & 1.0 & 32.1 \\
\addlinespace[1pt]
Table 4 (4): Set 2, constant error & 43.0 & 6.5 & 50.5 \\
\addlinespace[1pt]
Table 4 (5): Set 2, reliability predictors & 49.4 & 6.1 & 44.5 \\
\addlinespace[1pt]
Table 5 (2): AR(2) & 22.1 & 24.3 & 53.6 \\
\addlinespace[1pt]
Table 5 (6): AR(2), reliability predictors & 38.5 & 18.0 & 43.5 \\
\addlinespace[1pt]
Table 6 (3/4): FMM, common error & 31.5 & 19.6 & 48.8 \\
\addlinespace[1pt]
Table 6 (7/8): FMM, type intercepts & 27.5 & 18.7 & 53.9 \\
\addlinespace[1pt]
Table 7 (4): independent duration error* & 50.8 & 4.2 & 45.0 \\
\addlinespace[1pt]
Table 9 (1): duration AR(1)* & 56.4 & 6.0 & 37.6 \\
\addlinespace[1pt]
Table 9 (2): duration AR(2)* & 50.8 & 12.1 & 37.1 \\
\addlinespace[1pt]
\addlinespace
\multicolumn{4}{l}{\textit{Apparent entries}} \\
Table 1 (1): raw/no error & 0.0 & 8.5 & 91.5 \\
\addlinespace[1pt]
Table 1 (2): symmetric error & 66.0 & 1.0 & 33.0 \\
\addlinespace[1pt]
Table 4 (4): Set 2, constant error & 43.8 & 6.7 & 49.5 \\
\addlinespace[1pt]
Table 4 (5): Set 2, reliability predictors & 47.7 & 6.4 & 45.9 \\
\addlinespace[1pt]
Table 5 (2): AR(2) & 21.8 & 24.7 & 53.5 \\
\addlinespace[1pt]
Table 5 (6): AR(2), reliability predictors & 36.3 & 19.0 & 44.7 \\
\addlinespace[1pt]
Table 6 (3/4): FMM, common error & 31.0 & 20.2 & 48.7 \\
\addlinespace[1pt]
Table 6 (7/8): FMM, type intercepts & 27.4 & 19.2 & 53.4 \\
\addlinespace[1pt]
Table 7 (4): independent duration error* & 51.9 & 3.2 & 44.9 \\
\addlinespace[1pt]
Table 9 (1): duration AR(1)* & 61.7 & 6.5 & 31.8 \\
\addlinespace[1pt]
Table 9 (2): duration AR(2)* & 55.7 & 14.0 & 30.3 \\
\addlinespace[1pt]
\addlinespace
\multicolumn{4}{l}{\textit{Apparent exits}} \\
Table 1 (1): raw/no error & 0.0 & 8.3 & 91.7 \\
\addlinespace[1pt]
Table 1 (2): symmetric error & 67.9 & 1.0 & 31.1 \\
\addlinespace[1pt]
Table 4 (4): Set 2, constant error & 42.2 & 6.3 & 51.4 \\
\addlinespace[1pt]
Table 4 (5): Set 2, reliability predictors & 51.1 & 5.7 & 43.2 \\
\addlinespace[1pt]
Table 5 (2): AR(2) & 22.5 & 23.8 & 53.7 \\
\addlinespace[1pt]
Table 5 (6): AR(2), reliability predictors & 40.8 & 17.0 & 42.2 \\
\addlinespace[1pt]
Table 6 (3/4): FMM, common error & 32.1 & 19.0 & 48.9 \\
\addlinespace[1pt]
Table 6 (7/8): FMM, type intercepts & 27.5 & 18.1 & 54.4 \\
\addlinespace[1pt]
Table 7 (4): independent duration error* & 49.7 & 5.2 & 45.2 \\
\addlinespace[1pt]
Table 9 (1): duration AR(1)* & 50.9 & 5.5 & 43.6 \\
\addlinespace[1pt]
Table 9 (2): duration AR(2)* & 45.7 & 10.1 & 44.2 \\
\addlinespace[1pt]
\bottomrule
\end{tabular}\\[4pt]
{\footnotesize\textit{Notes.} Entries are survey-weighted percentages. Unstarred rows condition only on the two status reports forming an apparent transition and integrate out the subsequent report. Classification only means reported status changes but latent status does not. Genuine reversal means latent status changes and returns next period; genuine persistence means it remains in the new state. These outcomes are mutually exclusive and exhaustive; a latent change opposite to the reported direction is still a genuine change. Paired Table 6 references are the two types of one FMM, integrated over membership. *Table 7 and Table 9 use full-record posteriors including duration evidence, not status-pair-only conditioning. Table 9 pools apparent transitions 1--2 and 2--3 with their follow-up waves, using all four reports. Its rates are conditional on apparent switching, not the population status-error probability. Models use their own samples. Table 9 uses the current unrevised onset and employer-return assumptions; its uncertainty is pending.}
\end{table}
